\documentclass[10pt,a4paper]{amsart}
\usepackage[margin=19mm]{geometry}
\usepackage{amsmath,amssymb,mathtools}
\usepackage[scr=rsfs,cal=euler]{mathalfa}
\usepackage{xfrac}
\usepackage[unicode,hidelinks,bookmarks=false]{hyperref}
\newcommand{\ie}{i.e.\ }
\newcommand{\cf}{cf.\ }
\newcommand{\resp}{respectively\ }
\newcommand{\Subsection}{\subsection}
\providecommand{\nobreakdash}{\nobreak\hbox{-}}
\setlength{\emergencystretch}{2em}
\multlinegap0pt
\usepackage{booktabs}
\usepackage{nicefrac}
\usepackage{xcolor}
\newtheorem{corollary}{Corollary}
\title{Enhancing Diffusion Model Guidance through Calibration and Regularization}
\author{Seyed Alireza Javid, Amirhossein Bagheri,, Nuria Gonz\'alez-Prelcic}
\date{Source: arXiv:2511.05844v3 (CC BY 4.0)}
\begin{document}
\maketitle

\noindent\emph{Source and licence.} Enhancing Diffusion Model Guidance through Calibration and Regularization. Version arXiv:2511.05844v3 (CC BY 4.0), distributed by arXiv under the Creative Commons Attribution 4.0 International licence, http://creativecommons.org/licenses/by/4.0/, as recorded in the arXiv licence metadata for this exact version and shown on the abstract page https://arxiv.org/abs/2511.05844v3. Full licence notice: \texttt{licenses/arxiv-2511.05844-CC-BY-4.0.txt}.\par\smallskip

\noindent\textit{Editorial scope.} Counted unit: the article's corollary cor:rkl\_gradient (source0.tex lines 311--320). The article's complete original proof of it is delivered with it (source0.tex lines 1264--1279, one \texttt{proof} environment of 16 lines, which the article places away from the statement). No mathematics is written, repaired or re-derived: the statement and the proof are the article's own lines, in the article's own notation, and no retained line is substituted or re-flowed.  No further source-local context is needed: every cross-reference of every retained line resolves inside the two delivered spans.  Nothing was omitted from any retained span, so the package carries no omission record.  No retained line contains a citation, so the wrapper carries no bibliography and no bibliography entry.  No retained line contains a figure environment, an image-inclusion command or drawing code, so no image is delivered and the package declares no asset dependency.  Editorial formatting only, in addition: an AMS \texttt{amsart} preamble that restates the article's own declarations and macros used by the retained lines, namely the environments \texttt{corollary} on separate counters, together with the standard packages those lines need.  The retained environments print in this document as Corollary 1 in order of appearance, whereas the article prints the same items with the numbering of its own full document.  The complete native source of the article as distributed in the arXiv e-print for this exact version is bundled unchanged as \texttt{sources/188685/source0.tex}.
\begin{corollary}[Reverse KL gradient]
\label{cor:rkl_gradient}
The gradient of the reverse KL divergence regularized score is:{\fontsize{8.5}{10}\selectfont
\begin{equation}
\begin{aligned}
\nabla_x \mathcal{S}_{\text{RKL}}(x, y) =& \tau_1 \nabla_x f_y(x) - \tau_2 \sum_{i=1}^N p_{\tau_2}(i|x) \nabla_x f_i(x) \\&+ \alpha \sum_{i=1}^N q_y(i) \left[\nabla_x f_i(x) - \sum_{j=1}^N p(j|x) \nabla_x f_j(x)\right].
\end{aligned}
\label{eq:rkl_gradient}
\end{equation}}
\end{corollary}

\begin{proof}
Let $\ell_k(x)=\log b_k+\log f_k(x)$. For Gaussian components, $
\nabla_x \ell_k(x)=\nabla_x \log f_k(x)=\Sigma_k^{-1}(\mu_k-x)$. Additionally, since $\nabla_x \log f_k(x) = \frac{1}{f_k(x)} \nabla_x f_k(x) \rightarrow \nabla_x f_k(x) = f_k(x) \nabla_x \log f_k(x) $. Substituting these into Corollary~\ref{cor:rkl_gradient} gives
{\fontsize{8.5}{10}\selectfont
\begin{equation}
\begin{aligned}
\nabla_x \mathcal{S}_{\text{RKL}}(x, y)
&= \tau_1 f_y(x) \nabla_x \ell_y(x) - \tau_2 \sum_{k=1}^K p_{\tau_2}(k|x) f_k(x) \nabla_x \ell_k(x) \\
&\quad + \alpha \sum_{k=1}^K q_y(k) \Big(  f_k(x) \nabla_x \ell_k(x) - \sum_{j=1}^K  f_j(x)  w_j(x)\nabla_x \ell_j(x)\Big) \\
&= \sum_{k=1}^K  f_k(x) \Sigma_k^{-1}(\mu_k-x)\\ &\times\Big[\tau_1\mathbb{I}_{k=y}-\tau_2 p_{\tau_2}(k|x)+\alpha q_y(k)\Big] \\ &
     \;-\;\alpha\Big(\sum_{k=1}^K q_y(k)\Big)\sum_{j=1}^K f_j(x) w_j(x)\Sigma_j^{-1}(\mu_j-x).
\end{aligned}
\end{equation}}
Since $\sum_k q_y(k)=1$, collecting the coefficient of each $ f_k(x) \Sigma_k^{-1}(\mu_k-x)$ yields
$\Gamma_k(x,y,\alpha,\epsilon)=\tau_1\mathbb{I}_{k=y}-\tau_2 p_{\tau_2}(k|x)+\alpha(q_y(k)-w_k(x))$ as stated.
\end{proof}

\end{document}
